\chapter{API Deployment}\label{ch:deploy}

\textbf{Design.} Following the assignment, the \emph{model} goes to the cloud as an HTTP API and
the Flask web app runs on my machine and calls it. I chose \textbf{pathway C}: one Ubuntu VM on
\textbf{Azure} (Azure for Students: \$100 credit with a university e-mail, no credit card)
running a small Compose stack.

\textbf{API URL: \url{\APIURL}}

\begin{table}[H]\centering\small
\begin{tabularx}{\linewidth}{l X}
\toprule
Compose / local element & Azure equivalent used \\
\midrule
Host running Docker & VM \texttt{vm-23L-0764-A2-Ops}, Ubuntu 24.04 x64, size B2s (2 vCPU, 4\,GB) \\
Published ports 80/443 & Network Security Group inbound rules + \texttt{ufw} on the VM \\
\texttt{api} service & \file{api/Dockerfile}: gunicorn + Flask API + TensorFlow-CPU, model built in \\
HTTPS & Caddy reverse proxy with an automatic Let's Encrypt certificate for \texttt{<ip>.sslip.io} \\
Secrets (\file{.env}) & \texttt{API\_KEY} exported in the VM shell, never committed \\
Database & stays with the web app on my machine (Task 8 design) \\
\bottomrule
\end{tabularx}
\caption{Mapping the stack to Azure.}\end{table}

\textbf{Things researched first.}
\begin{itemize}
\item \emph{Free terms}: Azure for Students gives \$100 of credit valid for 12 months; a B2s VM
costs a few cents per hour, so a week of uptime is a small fraction of the credit. I stop the VM
when not needed.
\item \emph{Size}: TensorFlow plus the model needs roughly 1\,GB of RAM at start-up; the free
B1s (1\,GB) would be killed by the OOM killer, so I used B2s and also added 2\,GB swap.
\item \emph{x86 vs Arm}: x86-64 image, so the same \texttt{tensorflow-cpu} wheel as on my laptop and
in CI is used (Arm needs different wheels).
\item \emph{Two firewalls}: the cloud-level Network Security Group must allow 80/443 \emph{and}
the VM's own firewall (\texttt{ufw}) must too; either one alone blocks traffic.
\item \emph{Protecting the API}: the prediction endpoint requires a secret \texttt{X-API-Key}
header, compared in constant time (\texttt{hmac.compare\_digest}) so response timing does not
leak the key; HTTPS means the key is never sent in clear text; request bodies are capped at
6\,MB by Caddy and 5\,MB by the app; \texttt{/health} is public but reveals only the model version.
\end{itemize}

\section*{Step-by-step deployment}
\begin{enumerate}
\item \textbf{Create the VM} in the Azure portal: Ubuntu Server 24.04 LTS x64, size B2s, SSH key
authentication (download the \texttt{.pem}), inbound ports SSH, HTTP, HTTPS.
\shot{s7_vm_created}{VM \texttt{vm-23L-0764-A2-Ops} created (overview with public IP).}
\shot{s7_nsg_rules}{Network Security Group inbound rules: 22, 80, 443.}
\item \textbf{Connect and prepare}: \cmd{ssh -i key.pem azureuser@<ip>}, \cmd{git clone} the
repository, \cmd{bash deploy/vm_setup.sh}: installs Docker Engine and the compose plugin, adds the
user to the \texttt{docker} group, enables \texttt{ufw} allowing only SSH/80/443, and creates
2\,GB swap.
\shot{s7_ufw}{The VM's own firewall (\cmd{sudo ufw status}).}
\item \textbf{Ship the model}: \cmd{scp models/champion.keras models/champion.json ...:~/23L-0764_A2_Ops/models/}
(the champion exported from the registry; the file is git-ignored).
\item \textbf{Secrets and domain}: \cmd{export API_KEY=$(python3 -c "import secrets;print(secrets.token_urlsafe(32))")}
creates a 256-bit random key; \cmd{export API_DOMAIN=$(curl -s ifconfig.me | tr . -).sslip.io}
builds a hostname such as \texttt{20-55-1-9.sslip.io}, which the free sslip.io DNS resolves to
that IP, so Let's Encrypt can issue a certificate without buying a domain.
\item \textbf{Start}: \cmd{docker compose -f deploy/docker-compose.cloud.yml up -d --build}.
The API image is built with the model inside; Caddy starts only after the API is healthy, obtains
the certificate and proxies \texttt{https://API\_DOMAIN} to \texttt{api:8080}. Only Caddy publishes
ports; the API container is reachable only on the internal network.
\shot{s7_compose_up}{Cloud stack running: API healthy, Caddy up.}
\item \textbf{Verify from another network}: \texttt{/health} opened on my phone over mobile data.
\shot[0.5]{s7_health_phone}{\texttt{\APIURL/health} on a phone using mobile data (different network).}
\end{enumerate}

\section*{Testing with Postman}
The collection \file{deploy/VisionGuard.postman_collection.json} has variables \texttt{base\_url}
and \texttt{api\_key} and five requests with test scripts that assert the status code and the
response fields:
\begin{enumerate}
\item \texttt{GET /health} $\rightarrow$ 200 and the model name/version;
\item \texttt{POST /v1/predict} with an AI image (form-data field \texttt{image}) $\rightarrow$ 200, label, confidence, latency;
\item the same with a real photo;
\item without the \texttt{X-API-Key} header $\rightarrow$ 401;
\item with a \texttt{.txt} file $\rightarrow$ 415.
\end{enumerate}
\shot{s7_postman_health}{Postman: health.}
\shot{s7_postman_ai}{Postman: AI image $\rightarrow$ verdict JSON.}
\shot{s7_postman_real}{Postman: real photo $\rightarrow$ verdict JSON.}
\shot{s7_postman_401}{Postman: missing API key rejected (401).}
\shot{s7_postman_415}{Postman: wrong file type rejected (415).}
\shot{s7_postman_console}{Postman console: logged results of the run.}
\shot{s7_send_samples}{\file{scripts/send_samples.py} calling the cloud API with six test images.}

\yours{Copy one JSON response here (label, confidence, prob\_ai, model\_version, latency\_ms) and
the measured latency range, e.g.\ ``50--80\,ms model time, 300--500\,ms round trip from Lahore''.}
